% Generated draft: requires mathematical and source review.
\documentclass[11pt]{article}
\usepackage[margin=25mm]{geometry}
\usepackage{fontspec}
\setmainfont{Latin Modern Roman}
\usepackage{amsmath,amssymb,mathtools,unicode-math}
\setmathfont{Latin Modern Math}
\usepackage{enumitem,needspace,xurl,hyperref}
\setlength{\parindent}{0pt}
\setlength{\parskip}{5pt}
\setlength{\emergencystretch}{3em}
\newcommand{\N}{\mathbb{N}}
\newcommand{\Z}{\mathbb{Z}}
\newcommand{\Q}{\mathbb{Q}}
\newcommand{\R}{\mathbb{R}}
\newcommand{\C}{\mathbb{C}}
\newcommand{\F}{\mathbb{F}}
\newcommand{\A}{\mathbb{A}}
\newcommand{\PP}{\mathbb{P}}
\newcommand{\E}{\mathbb{E}}
\newcommand{\eps}{\varepsilon}
\newcommand{\dd}{\,\mathrm{d}}
\begin{document}
\section*{UnsolvedMath ID: 30005465}
\textbf{Source problem code:} catalogue entry
\subsection{Definitions and mathematical statement}
The supplied record asks for a characterization of the intersection of the sum-of-squares and SONC cones. It does not explicitly define their ambient space or parameter range. The following gives a conventional real-polynomial interpretation; these conventions require confirmation against the cited report.

For positive integers \(n,d\), let
\[
V_{n,2d}=\{p\in\R[x_1,\ldots,x_n]:\deg p\leq 2d\},
\qquad
\Sigma_{n,2d}=\left\{\sum_{j=1}^{m}q_j^2:m\in\Z_{\geq0},\ q_j\in\R[x_1,\ldots,x_n],\ \deg q_j\leq d\right\}.
\]
Empty sums are understood to be zero. For \(\alpha\in\Z_{\geq0}^{n}\), write \(x^\alpha=\prod_{k=1}^{n}x_k^{\alpha_k}\) and \(|\alpha|=\sum_{k=1}^{n}\alpha_k\).

Under the same convention, a circuit polynomial of degree at most \(2d\) has the form
\[
f(x)=\sum_{i=0}^{r}a_i x^{\alpha_i}+b x^\beta,
\]
where \(1\leq r\leq n\), the vectors \(\alpha_0,\ldots,\alpha_r\in(2\Z_{\geq0})^n\) are affinely independent, \(a_i>0\), \(b\in\R\setminus\{0\}\), and
\[
\beta\in\Z_{\geq0}^{n}\cap\operatorname{relint}\operatorname{conv}\{\alpha_0,\ldots,\alpha_r\},
\qquad |\alpha_i|\leq2d,\quad |\beta|\leq2d.
\]
Here \(\operatorname{relint}\) denotes interior relative to the affine hull. Such a polynomial is nonnegative if \(f(x)\geq0\) for every \(x\in\R^n\). Let \(C_{n,2d}\) be the cone of finite sums of these nonnegative circuit polynomials and nonnegative monomials \(c x^{2\gamma}\), where \(c\geq0\), \(\gamma\in\Z_{\geq0}^{n}\), and \(|\gamma|\leq d\). This is the conventional sum-of-nonnegative-circuit-polynomials (SONC) interpretation used in this draft.

\textbf{Problem.} Characterize
\[
\Sigma_{n,2d}\cap C_{n,2d};
\]
that is, give necessary and sufficient conditions describing the polynomials belonging to both cones, for the parameters intended in the source.

The supplied excerpt does not settle whether the intended setting is all polynomials of total degree at most \(2d\), homogeneous forms of degree \(2d\), or another convention. It also does not specify parameter restrictions, treatment of degree zero, or the desired form of a characterization. The definitions above are therefore a provisional reconstruction rather than definitions quoted from the report.
\subsection{Short English statement}
Describe exactly which polynomials are both sums of squares and sums of nonnegative circuit polynomials. The draft uses conventional real-polynomial definitions, whose agreement with the source requires review. The supplied dated curation labels the general characterization open.
\subsection{Sources}
\begin{itemize}
\item Supplied local registry record: permanent UnsolvedMath ID \(30005465\), title \emph{Intersection of Sum-of-Squares and SONC Cones}, exact source problem code \texttt{catalogue entry}, in the Oberwolfach Reports Open Problems collection. The recorded proposed year is 2023; no individual proposer is identified.
\item \emph{Real Algebraic Geometry with a View toward Koopman Operator Methods} (2023), Oberwolfach Reports, pp.\ 741--811, DOI 10.4171/owr/2023/14. Supplied report provenance and literature source: \url{https://doi.org/10.4171/owr/2023/14}.
\item Dressler--Kuhlmann--Schick, \emph{Advances in Geometry}, DOI 10.1515/advgeom-2024-0031. Supplied literature source: \url{https://doi.org/10.1515/advgeom-2024-0031}. No article title or further bibliographic details are supplied.
\item Literature assessment dated exactly \texttt{2026-08-22}: the supplied curation labels this problem open and states that surrounding SOS+SONC equality questions were resolved, while the report and subsequent primary paper leave a general characterization of the intersection open. This is supplied dated curation, not independently verified literature assessment; its status and cited claims require independent checking.
\item The registry labels the statement \texttt{corrected\_verified} and reports that the official report was checked directly and extraction artifacts were removed. This verification claim is supplied metadata; no independent source inspection is asserted here.
\end{itemize}
\end{document}
